\section{A constructive amplifier on well-founded trees}
\label{amp:construction}

The amplification mechanism is independent of the benchmark enumerator.  Its
input is an explicitly defined function $b:\mathbb N\to\mathbb N$ satisfying
$b(n)\geq n+1$.  Its output is a closed natural number $N_b$.  The construction
has two complementary features.  Countably branching well-founded trees and
operators on higher function spaces provide a substantial supply of recursive
structure.  A separate feedback mechanism supplies a quantitative growth
certificate that does not require assigning ordinal names to those trees.
The result is an exact definition together with an independently checkable
lower bound.

Throughout this section, iteration means ordinary finite iteration:
\[
 f^{[0]}(x)=x,
 \qquad f^{[r+1]}(x)=f^{[r]}(f(x)).
\]
This convention exposes the first application of $f$.  It defines the same
iterate as the familiar recurrence $f^{[r+1]}(x)=f(f^{[r]}(x))$.
The brackets distinguish iteration from exponentiation.

\subsection{Trees, evaluation, and substitution}
\label{amp:trees}

Let $\mathcal T$ be the inductively generated type with constructors
\[
 0_{\mathcal T},\qquad
 S(a)\quad(a\in\mathcal T),\qquad
 L(f)\quad(f:\mathbb N\to\mathcal T).
\]
The third constructor packages a countable family of immediate subtrees.
The word ``inductively'' is essential: these are well-founded trees, and their
eliminator permits structural recursion on a tree, including recursive calls
to any member $f(r)$ of a family appearing at a limit node.  They are not
coinductive trees, and a cyclic equation does not introduce a new inhabitant.
We use $0_{\mathcal T}$ rather than the numeral $0$ when its type might be
ambiguous.

For a base $b:\mathbb N\to\mathbb N$, define
\begin{align}
 F_b(0_{\mathcal T};n)&=b(n),\label{amp:eq:fastzero}\\
 F_b(S(a);n)&=\bigl(F_b(a;\mathord\cdot)\bigr)^{[n+1]}(n),
 \label{amp:eq:fastsucc}\\
 F_b(L(f);n)&=F_b(f(n);n).
 \label{amp:eq:fastlimit}
\end{align}
Thus a successor node turns the function supplied by its child into a finite
iterator.  A limit node selects its $n$th child.  The argument $n$ is allowed
to change inside a successor iteration; consequently the evaluator genuinely
uses the functions associated with subtrees, rather than merely counting
nodes along one fixed path.

For $a,t\in\mathcal T$, define the substitution $a\star t$ by replacing every
zero leaf of $t$ with $a$:
\begin{align*}
 a\star0_{\mathcal T}&=a,\\
 a\star S(t)&=S(a\star t),\\
 a\star L(f)&=L(r\mapsto a\star f(r)).
\end{align*}
The orientation matters: the first argument is inserted into the second.
The Lean implementation calls this operation \texttt{graft}.

The evaluator returns a function of a natural argument at each tree.  Its
successor equation can therefore reuse the child function with increasingly
large arguments, although the recursive call defining that function is
still to a smaller tree.  This distinction is useful when reading the
termination argument: the natural arguments to child functions need not
decrease.  Structural recursion decreases the tree argument, while the
separate iteration operation decreases its finite iteration counter.
The construction never asks the evaluator to recurse on a tree that it has
just generated from itself.

\begin{lemma}[Inflationarity and the base bound]
\label{amp:lem:base}
If $b(n)\geq n$ for every $n$, then for every $a\in\mathcal T$ and every $n$,
\[
 n\leq b(n)\leq F_b(a;n).
\]
If instead $b(n)\geq n+1$ for every $n$, then
$F_b(a;n)\geq n+1$ for every $a,n$.
No monotonicity hypothesis on $b$ is required.
\end{lemma}

\begin{proof}
First observe that $g(x)\geq x$ implies $g^{[r]}(x)\geq x$ for all $r,x$;
this follows by induction on $r$.  Now use structural induction on $a$.
The zero case is equality.  In the successor case, put
$g(x)=F_b(a;x)$.  The induction hypothesis gives $g(x)\geq b(x)\geq x$.
Therefore
\[
 F_b(S(a);n)=g^{[n]}(g(n))\geq g(n)\geq b(n).
\]
At a limit node, use the induction hypothesis for the selected child $f(n)$.
This proves the weak statement.  The strict statement follows by first
weakening $b(n)\geq n+1$ to $b(n)\geq n$ and then retaining the stronger first
inequality in $F_b(a;n)\geq b(n)\geq n+1$.
\end{proof}

The absence of a monotonicity assumption is more than a simplification of
the statement.  An inflationary function may take a smaller value at a
larger argument; nevertheless its own orbit is nondecreasing, because each
successive orbit value is at least its immediate predecessor.  The proof
uses precisely this orbit property.  It does not replace $b(x)$ by $b(y)$
merely because $x\leq y$.  This makes the interface applicable to bases
assembled from finite searches whose dependence on their search parameter
has not separately been proved monotone.

\begin{theorem}[Substitution law]
\label{amp:thm:graft}
For every base $b$, without an inflationarity assumption, every $a,t$, and
every $n$,
\[
 F_b(a\star t;n)=F_{F_b(a;\mathord\cdot)}(t;n).
\]
\end{theorem}

\begin{proof}
Proceed by structural induction on $t$.  At a zero leaf both sides are
$F_b(a;n)$.  For a successor node, the induction hypothesis identifies the
two child functions at every natural argument.  Pointwise equal functions
have pointwise equal finite iterates: prove this auxiliary assertion by
induction on the iteration count, replacing the first application and then
the remaining iterate.  Applying it with count $n+1$ gives the successor
case.  At a limit node, both evaluators select the child indexed by $n$, and
the induction hypothesis for that child applies directly.  The proof uses
pointwise equalities throughout.  In particular, the Lean proof does not
invoke function extensionality.
\end{proof}

\begin{proposition}
\label{amp:cor:graftlower}
If $b(n)\geq n$ everywhere, then
\[
 F_b(a\star t;n)\geq F_b(a;n)
 \qquad(a,t\in\mathcal T).
\]
\end{proposition}

\begin{proof}
Lemma~\ref{amp:lem:base} says that $F_b(a;\mathord\cdot)$ is an inflationary
base.  Apply the same lemma to that base and to $t$, and then use
Theorem~\ref{amp:thm:graft}.
\end{proof}

The substitution law is the main algebraic safeguard in the construction.
It permits an elaborate tree to be attached to an existing tree while
preserving a proved evaluation bound.  The later operators need not satisfy
an unproved comparison principle for ordinal notations.

More specifically, $t$ acts as a program transforming the initial unary
function supplied at its zero leaves.  Substituting $a$ into those leaves
replaces that initial function by the already constructed function
$F_b(a;\mathord\cdot)$.  The theorem identifies these two operations exactly,
not merely up to a lower bound.  The tree $t$ may depend on $a$, on an
evaluation of $a$, or on other previously defined data.  Once those data
have been supplied, the same substitution theorem applies to the resulting
pair of trees.  That flexibility is what permits the later feedback
construction without weakening its proof obligation.

\begin{proposition}[The zero argument and successor gain]
\label{amp:prop:gain}
For every base $b$ and tree $a$,
\[
 F_b(a;0)=b(0).
\]
If $b(n)\geq n+1$ everywhere, then
\[
 F_b(S(a);n)\geq F_b(a;n)+n.
\]
\end{proposition}

\begin{proof}
The first assertion follows by tree induction.  A successor performs exactly
one child application at argument zero, and a limit selects its zeroth
child, so both cases reduce to the corresponding induction hypothesis.
For the second assertion, an induction on $r$ gives
$g^{[r]}(x)\geq x+r$ whenever $g(x)\geq x+1$ for all $x$.
Apply this to $g=F_b(a;\mathord\cdot)$, which has the required property by
Lemma~\ref{amp:lem:base}, and calculate
\[
 F_b(S(a);n)=g^{[n]}(g(n))\geq g(n)+n.
\]
The gain vanishes at $n=0$, exactly as the first assertion requires.
\end{proof}

For a small concrete calibration, take $b(n)=n+1$ and finite successor trees
$a_0=0_{\mathcal T}$, $a_{r+1}=S(a_r)$.  Directly from the definition,
\[
 F_b(a_0;n)=n+1,\qquad F_b(a_1;n)=2n+1,
\]
and
\[
 F_b(a_2;n)=2^{n+1}(n+1)-1.
\]
Indeed, the $r$fold iterate of $x\mapsto2x+1$ at $n$ is
$2^r(n+1)-1$, as an induction on $r$ verifies.  Substituting $r=n+1$ gives
the displayed formula.  These elementary calculations concern the
successor-only fragment; they do not supply an ordinal interpretation of
the general limit constructor.
The next tree iterates the last function $n+1$ times.  These identities
explain the basic acceleration before countable branching or higher
function spaces enter the definition.  They also show why the evaluation
argument is part of the construction, rather than an incidental input.

\subsection{Explicit tree operators}
\label{amp:operators}

We now specify the entire additional tree payload.  The following definitions
are syntactic operations on $\mathcal T$.  Their resemblance to familiar
ordinal constructions guides the design, but no identification with a
Veblen ordinal or an ordinal collapsing function is assumed in any proof.

For a tree transformation $f:\mathcal T\to\mathcal T$, set
\[
 O_f(a)=L\bigl(r\mapsto f^{[r+1]}(a)\bigr).
\]
This \emph{orbit-limit constructor} packages successively longer finite
orbits.  The name does not assert that it produces a fixed point of an
arbitrary $f$.

Define a tree transformation $W$ by
\begin{align*}
 W(0_{\mathcal T})&=S(0_{\mathcal T}),\\
 W(S(a))&=L\bigl(r\mapsto(x\mapsto W(a)\star x)^{[r]}(0_{\mathcal T})\bigr),\\
 W(L(f))&=L(r\mapsto W(f(r))).
\end{align*}
The implementation names this transformation \texttt{omegaPower}.  In
particular, its successor clause iterates substitution into a zero tree;
it is not an appeal to an externally supplied ordinal exponentiation
operation.

For a supplied transformation $c:\mathcal T\to\mathcal T$, define
\begin{align*}
 E_c(0_{\mathcal T})&=c(0_{\mathcal T}),\\
 E_c(S(a))&=c(S(E_c(a))),\\
 E_c(L(f))&=L(r\mapsto E_c(f(r))).
\end{align*}
This is an enumeration constructor generated by $c$.  For a countable family
$\boldsymbol f=(f_0,f_1,\ldots)$ of tree transformations, also define finite
prefix applications and repeated sweeps:
\begin{align*}
 P_{\boldsymbol f}(0,a)&=a,&
 P_{\boldsymbol f}(r+1,a)&=f_r(P_{\boldsymbol f}(r,a)),\\
 Q_{\boldsymbol f}(0,a)&=a,&
 Q_{\boldsymbol f}(r+1,a)&=
 P_{\boldsymbol f}(r+1,Q_{\boldsymbol f}(r,a)).
\end{align*}
The common-limit constructor is
\[
 C_{\boldsymbol f}(a)=L(r\mapsto Q_{\boldsymbol f}(r+1,a)).
\]
Repeated sweeps deliberately revisit earlier members of the family.  This
choice is more structured than applying each $f_r$ once and then discarding
it.  However, our correctness argument only needs the explicit recurrence;
it does not require a fixed-point theorem for $C_{\boldsymbol f}$.

Define a two-argument transformation $V$ by recursion on its first tree:
\begin{align*}
 V(0_{\mathcal T},a)&=W(a),\\
 V(S(t),a)&=E_{O_{V(t,\mathord\cdot)}}(a),\\
 V(L(f),a)&=
 E_{C_{(V(f(r),\mathord\cdot))_{r\in\mathbb N}}}(a).
\end{align*}
Here $O_{V(t,\mathord\cdot)}$ is the transformation
$x\mapsto O_{V(t,\mathord\cdot)}(x)$, and similarly for the common-limit
constructor in the last line.  The recursive calls $V(t,\mathord\cdot)$ and
$V(f(r),\mathord\cdot)$ involve immediate subtrees of the first argument.
Finally, set
\[
 \gamma=O_{a\mapsto V(a,0_{\mathcal T})}(0_{\mathcal T}).
\]
The symbol $\gamma$ denotes this particular tree.  It is not a claim that
the tree has been proved to represent $\Gamma_0$.

These definitions produce one explicit transformation rather than leaving
a black-box ``large ordinal'' in the candidate.  Every tree used below is
built from the three constructors and the recurrences just displayed.
Their role in the size certificate is equally explicit: a resulting tree
is placed in the second argument of substitution, where
Proposition~\ref{amp:cor:graftlower} applies regardless of its further behavior.

The dependency order also explains totality of this part of the definition.
Substitution is available before $W$ is defined.  The orbit, enumeration,
prefix, and sweep operations are available before $V$ is defined.  In the
successor clause of $V$, the transformation associated with the smaller
first tree has already been supplied by structural recursion; enumeration
then acts on the second tree using that transformation as a parameter.
The limit clause receives the analogous family of transformations for the
children of its first tree.  There is no recursive call to $V$ on a newly
constructed first argument in these clauses.

\subsection{A dynamically chosen tower of function spaces}
\label{amp:levels}

The next layer iterates operators at successively higher function types.
Define
\[
 \mathcal L_0=\mathcal T,
 \qquad \mathcal L_{k+1}=\mathcal L_k\to\mathcal L_k.
\]
Thus an element of $\mathcal L_{k+2}$ transforms endomorphisms of
$\mathcal L_k$.  Each index is a natural number, and the family itself is
defined by natural-number recursion.  No operation indexed by an infinite
natural number is being introduced.

Successor and countable branching extend pointwise through these spaces.
For $x\in\mathcal L_k$ and $f:\mathbb N\to\mathcal L_k$, define
\begin{align*}
 s_0(x)&=S(x),&
 s_{k+1}(g)(x)&=s_k(g(x)),\\
 \ell_0(f)&=L(f),&
 \ell_{k+1}(f)(x)&=\ell_k(r\mapsto f(r)(x)).
\end{align*}
If $g\in\mathcal L_{k+1}$, put
\[
 \Omega_k(g)(x)=
 \ell_k\bigl(r\mapsto g^{[r+1]}(s_k(x))\bigr).
\]
The successor $s_k(x)$ is inserted before the finite orbit.  This detail is
part of the exact definition and is retained at every function level.
The operator $\Omega_k$ belongs to $\mathcal L_{k+2}$.

Let
\[
 B_0=W,\qquad B_{k+1}=\Omega_k.
\]
In both cases $B_k\in\mathcal L_{k+1}$.  A lowering operation converts the
result back into a transformation of trees:
\[
 D_0(f)=f,\qquad
 D_{k+1}(f)=D_k(f(B_k)).
\]
The types match at each step: if $f\in\mathcal L_{k+2}$, then
$f(B_k)\in\mathcal L_{k+1}$.  Define the seed family
\[
 \Sigma_k=D_k(B_k)(\gamma).
\]
This family is defined for every natural $k$.  Later, $k$ will itself be an
earlier evaluation score plus $64$.  The dependent type family therefore
provides a compact description of a finite function-space tower whose
height is determined by the construction, rather than fixed by a long
list of manually written arrow types.

The first three seeds make the difference between levels concrete:
\begin{align*}
 \Sigma_0&=W(\gamma),\\
 \Sigma_1&=\Omega_0(W)(\gamma)
   =L(r\mapsto W^{[r+1]}(S(\gamma))),\\
 \Sigma_2&=\Omega_1(\Omega_0)(W)(\gamma).
\end{align*}
In the last expression, $\Omega_0$ is itself the function being iterated by
the next operator.  Expanded one step further, it is
\[
 \Sigma_2=L\bigl(r\mapsto
   [\Omega_0^{[r+1]}(x\mapsto S(W(x)))](\gamma)\bigr).
\]
Thus the second increase of level is not just another application of $W$.
It iterates the operation that constructs an orbit-limit transformation.
Higher levels continue this same typed pattern.  The recurrence for $D_k$
ensures that every such operator is eventually applied to appropriate
arguments and produces an ordinary member of $\mathcal T$.

\subsection{Feedback and the final number}
\label{amp:feedback}

Fix the evaluation argument $64$, and abbreviate
\[
 \sigma_b(a)=F_b(a;64),\qquad
 t_b(a)=\Sigma_{\sigma_b(a)+64}.
\]
The elementary advancement operation is
\begin{equation}
 A_b(a)=S(a)\star
 V\bigl(a\star t_b(a),\ t_b(a)\star a\bigr).
 \label{amp:eq:advance}
\end{equation}
Both orders of substitution occur in the payload.  The outer substitution
has first argument $S(a)$; that choice makes the successor gain available
before any properties of the payload are investigated.

Let $T$ be the tower of twos with the following exact height convention:
\[
 T(0)=1,\qquad T(r+1)=2^{T(r)}.
\]
One complete stage performs a number of advancement operations given by a
tower whose height is the current score plus one:
\begin{equation}
 G_b(a)=A_b^{[T(\sigma_b(a)+1)]}(a).
 \label{amp:eq:stage}
\end{equation}
Starting with $a_0=\Sigma_{64}$, define
\[
 a_{r+1}=G_b(a_r),\qquad s_b=\sigma_b(a_0).
\]
The proposed output of the amplifier is
\begin{equation}
 \boxed{N_b=\sigma_b(a_{s_b})=F_b(a_{s_b};64).}
 \label{amp:eq:candidate}
\end{equation}
The actual score of the initial tree is therefore also the number of
complete feedback stages.  At each stage, scores additionally control both
a tower height and the level of a function-space construction.

There is no circular definition here.  First define tree evaluation and
the tree operators.  Then define $A_b$ using the already available evaluator
on its input tree.  Finite iteration defines $G_b$, natural-number recursion
defines $a_r$, and the already defined initial score supplies the final
finite index.  During iteration, successive trees can become more complex;
termination of that iteration is nevertheless justified by its decreasing
natural-number counter.  A defining term need not reduce its enormous
numeric arguments during type checking.

There are two distinct uses of a score inside a stage.  Its total iteration
count $T(\sigma_b(a)+1)$ is fixed from the input tree before that stage
begins.  Within each advancement, the current tree determines a new
function-space index for $t_b$.  Recomputing that index does not reset the
iteration counter.  Likewise, the total number $s_b$ of complete stages is
fixed by the initial seed.  These distinctions rule out an interpretation
as an open-ended loop whose stopping condition continually moves away.
Every iteration count is a finite natural number supplied to an ordinary
recursor.

\subsection{A quantitative certificate independent of calibration}
\label{amp:bounds}

Assume for the remainder of this section that $b(n)\geq n+1$ everywhere.
Equation~\eqref{amp:eq:advance}, Proposition~\ref{amp:cor:graftlower}, and
Proposition~\ref{amp:prop:gain} give, for every evaluation argument $n$,
\begin{equation}
 F_b(A_b(a);n)\geq F_b(S(a);n)\geq F_b(a;n)+n.
 \label{amp:eq:advancegain}
\end{equation}
The payload in $A_b$ uses the score at $64$, but the displayed comparison
holds at every $n$.  Induction on the number of advancements now yields
\[
 F_b(A_b^{[r]}(a);n)\geq F_b(a;n)+rn.
\]
Writing $m=\sigma_b(a)$ and applying this inequality to one complete stage
gives the sharper intermediate estimate
\begin{equation}
 \sigma_b(G_b(a))\geq m+64T(m+1).
 \label{amp:eq:rawstage}
\end{equation}

The elementary inequalities $2^q\geq q+1$ imply by induction that
$T(r)\geq r+1$.  They also imply that $T$ is increasing.  Consequently
\eqref{amp:eq:rawstage} has two useful consequences:
\begin{align}
 \sigma_b(G_b(a))&\geq T(m+1),
 \label{amp:eq:stage-tower}\\
 \sigma_b(G_b(a))+1&\geq65(m+1).
 \label{amp:eq:stage-scale}
\end{align}
The first follows because $m\geq0$ and $64\geq1$.  For the second, use
$T(m+1)\geq m+1$ to obtain
$m+64T(m+1)+1\geq m+64(m+1)+1=65(m+1)$.

Define a purely numerical comparison function by
\[
 \Theta(s,0)=s,\qquad
 \Theta(s,r+1)=T(\Theta(s,r)+1).
\]
This repeatedly replaces the previous number by a tower of twos whose
height is that number plus one.  Its recursion is useful for communicating
a guaranteed lower bound: it requires none of the tree operators or higher
function spaces.

For example, the first two numerical feedback values are
\[
 \Theta(s,1)=T(s+1),\qquad
 \Theta(s,2)=T(T(s+1)+1).
\]
The second expression uses the entire first tower value as part of a new
tower height.  The final bound uses $s$ such replacements.  It is therefore
substantially more explicit than saying only that the candidate contains
many iterations: it gives a familiar recursively defined natural function
that the actual candidate is proved to exceed.

\begin{theorem}[Certified feedback growth]
\label{amp:thm:quantitative}
For every $r\in\mathbb N$,
\begin{align*}
 \sigma_b(a_r)&\geq\Theta(s_b,r),\\
 \sigma_b(a_r)+1&\geq65^r(s_b+1).
\end{align*}
In particular,
\[
 \boxed{N_b\geq\Theta(s_b,s_b)},\qquad
 \boxed{N_b+1\geq65^{s_b}(s_b+1)}.
\]
\end{theorem}

\begin{proof}
Both claims are proved by induction on $r$.  At $r=0$, the first is equality
by the definition of $s_b$, and the second uses $65^0=1$.  Suppose the first
claim holds at $r$.  Monotonicity of $T$, followed by
\eqref{amp:eq:stage-tower}, gives
\[
 \Theta(s_b,r+1)
 =T(\Theta(s_b,r)+1)
 \leq T(\sigma_b(a_r)+1)
 \leq\sigma_b(a_{r+1}).
\]
For the second claim, multiply its induction hypothesis by $65$ and apply
\eqref{amp:eq:stage-scale}.  Associativity of multiplication and the
successor equation for powers give the required factor $65^{r+1}$.
Finally specialize both inequalities to $r=s_b$.
\end{proof}

\begin{proposition}[Transferring a benchmark into the bound]
\label{amp:prop:transfer}
The function $\Theta$ is nondecreasing in each argument.  Hence, whenever
$c\leq s_b$,
\[
 \Theta(c,c)\leq N_b.
\]
Also $b(64)\leq s_b$ and $b(64)\leq N_b$.
\end{proposition}

\begin{proof}
Monotonicity in the first argument follows by induction on the second,
using monotonicity of $T$.  Since $T(q+1)\geq q+2$, each step of the
second argument increases the value, proving its monotonicity as well.
Thus $c\leq s_b$ gives
$\Theta(c,c)\leq\Theta(s_b,c)\leq\Theta(s_b,s_b)\leq N_b$.
The two base comparisons are instances of Lemma~\ref{amp:lem:base}, first
for $a_0$ and then for $a_{s_b}$.
\end{proof}

The proof distinguishes a guaranteed quantitative margin from a proposed
ordinal interpretation.  The entire payload appears in the concrete
definition of $N_b$, but the displayed lower bounds remain valid even if
its ordinal calibration is left unresolved.  They follow from structural
recursion, substitution, inflationarity, and explicit finite iteration.
This makes the benchmark comparison robust while leaving room for further
analysis of how much additional growth the higher operators contribute.

There is a useful modular consequence of the argument.  Suppose a later
investigation supplies a different explicit total tree transformation
$R_b(a)$.  Replacing the payload in~\eqref{amp:eq:advance} by $R_b(a)$ while
retaining the outer form $S(a)\star R_b(a)$ preserves
\eqref{amp:eq:advancegain}.  The induction giving the raw stage estimate and
both final bounds then applies without changing its reasoning.  The
substitution interface therefore isolates a stable mathematical obligation
from the choice of additional tree-generating operators.  Such a replacement
would still require its own well-typed definition and resource checks; the
present candidate is the fully specified payload above, rather than an
unspecified member of that wider family.

Conversely, the numerical bound alone does not identify the exact value of
$N_b$.  The raw stage inequality already discards most information about
the generated payload, and the tower comparison discards the additional
factor $64$ and the previous score.  These are deliberate weakenings in a
proof of a lower bound.  They make the certificate easy to state and
transport to a benchmark, while the exact recurrences retain all the
additional structure in the proposed number.
